% Agent context and source locators: ../../../internal/EXPANSION_NOTES.md
\section{Preliminary Work and Research Readiness}\label{sec:preliminary-work-and-research-readiness}

\subsection{Master's work: an existing end-to-end reduction example}\label{subsec:master-s-work-an-existing-end-to-end-reduction-example}

\begin{itemize}[leftmargin=*]
\item The completed master\textquotesingle{}s work establishes feasibility for an end-to-end CFD-to-1D reduction process. The next experimental demonstration will test data interpretation and temperature reproduction in the ACU or VCU setting.\cite{fierro2026cfdoneD}
\item The TAMU study relates CFD fields and sensor meanings to matched mass, energy, pressure, and temperature quantities, with explicitly defined averaging windows.\cite{fierro2026cfdoneD}
\item The steady model includes energy transport around the loop, temperature periodicity, a buoyancy--resistance balance, separate thermal paths, and distinct fluid and wall observations.
\item Assess temperature versus loop location with the case population, supplied inputs, fitted choices, and sensor policy stated explicitly.
\item The master\textquotesingle{}s assessment is retrospective: nominal CFD cases informed development and the heater-power cases had been encountered previously. These results support feasibility and lessons for assessment design; they do not establish independent ACU or VCU validation.
\end{itemize}

\subsection{Quantitative evidence from the completed study}\label{subsec:quantitative-preliminary-evidence}
\begin{itemize}[leftmargin=*]
\item The completed study assessed four nominal CFD cases and six previously encountered heater-power cases. The reported balanced temperature metric combines the fluid-probe and wall-probe comparisons under the thesis's declared sensor-inclusion policy.\cite{fierro2026cfdoneD}
\begin{center}
\small
\begin{tabular}{@{}p{3.05in}rr@{}}
\toprule
\textbf{Model setting} & \multicolumn{2}{c}{\textbf{Balanced temperature RMSE (K)}}\\
 & \textbf{Four nominal cases} & \textbf{Six power cases}\\
\midrule
Global thermal-compensation setting & 9.280 & 12.853\\
Localized thermal setting & 4.957 & 5.989\\
Localized, flow-calibrated setting & 7.090 & 5.970\\
\bottomrule
\end{tabular}
\end{center}
\item The separate localized, flow-calibrated setting achieved CFD mass-flow mean absolute percentage errors of 0.967\% and 1.310\% for those two case groups. Its hydraulic setting was selected using known nominal CFD flow; the six power cases had also been seen previously. These are development and retrospective assessment results.\cite{fierro2026cfdoneD}
\item The lower temperature errors did not establish the correct heat-removal distribution. Among the three compared heat sinks, the localized model assigned 6.9--7.3\% to the cooler, compared with 57.4--60.8\% in CFD, and 56.1--57.0\% to junction regions, compared with 15.3--16.5\% in CFD.
\begin{itemize}
\item These fractions refer to the explicitly compared sinks, not an unrestricted fraction of all generated power. The fitted residual thermal conductance is an effective model term rather than a measured hardware property.
\item The result motivates a PhD assessment that scores measured temperatures, circulation evidence, and regional energy accounting separately. A single aggregate temperature score can conceal compensating heat-path errors.
\end{itemize}
\item The primary temperature metric omitted a declared subset of wall channels; restoring some of those channels increased the earlier global-compensation errors from roughly 9--10 K to roughly 14--15 K. A corresponding restored-channel grade was not reported for the localized model. Sensor-policy sensitivity is therefore part of the evidence, not an optional presentation choice.
\end{itemize}

\subsection{Lessons that motivate the PhD}\label{subsec:lessons-that-motivate-the-phd}
\begin{itemize}[leftmargin=*]
\item Temperature and circulation objectives can favor different fitted settings; show both rather than selecting a temperature score as evidence of correct flow.
\item A lower temperature error can coexist with incorrect allocation of cooler and junction heat removal. Carry local energy budgets alongside temperature agreement.
\item The primary wall-temperature aggregate used a declared channel-selection policy. Restoring previously omitted channels changes the earlier models\textquotesingle{} assessment; channel exclusion therefore requires explicit reporting and sensitivity analysis.
\item Use hydraulic/\allowbreak{}thermal development and observation meaning to choose candidate repairs; a fitted residual loss or resistance is not automatically a measured hardware property.
\item The tested two-stream heat exchange did not improve temperatures, and the fixed-inlet alternative violated mixing. Recirculation in TAMU motivates a bounded question; it does not establish ACU/\allowbreak{}VCU recirculation or a transferable coefficient.
\item PhD step beyond this work: assess a reproducible reduction/\allowbreak{}closure/\allowbreak{}observation procedure using qualified experiment and CFD evidence, with uncertainty and unused assessment conditions. Reproducing a temperature trace is an initial demonstration; an additional method benefit or supported limit is the research contribution.
\end{itemize}

\subsection{TRIDENT 1D groundwork and evidence status}\label{subsec:trident-1d-groundwork-and-evidence-status}
\begin{itemize}[leftmargin=*]
\item \researchlabel{ambient-and-startup-comparisons}\textbf{Ambient and startup-conditioned temperature comparisons.} Preliminary studies examine alternative ambient and initial-state assumptions using observed-versus-predicted temperature traces and channel residuals.
\begin{itemize}
\item Keep other-run learning, previously inspected windows and supplied later boundaries explicit; mixed performance is not a universal winner.
\end{itemize}

\item The physical model includes segmented fluid, pipe, and optional insulation states; zero, prescribed, or coupled steady flow; and conservation outputs. Synthetic examples demonstrate numerical behavior rather than facility prediction.
\item Ambient-only, startup-assumption, and startup-conditioned comparisons provide groundwork for assessing how supplied initial and boundary information affects temperature reproduction.
\end{itemize}

\subsection{Next feasibility demonstration: understand data and reproduce temperatures}\label{subsec:next-feasibility-demonstration-understand-data-and-reproduce-temperatures}

\begin{itemize}[leftmargin=*]
\item Assemble a documented experimental case with facility configuration, timestamps, units, channel mapping, calibration, heating, cooling and ambient histories, and known data gaps.
\item For each temperature, state whether it is a primary-fluid probe, wall/\allowbreak{}contact sensor or fiber measurement and specify the corresponding model observation. Separate measured values, averaged values and inferred bulk states.\cite{jcgm2020measurementmodels}
\item First reproduce an accepted development case using declared inputs and fitted quantities; report per-sensor residuals, spatial patterns, energy closure and uncertainty. Call this reproduction or calibration when target data inform the model.
\item Then freeze those choices and compare against a separate eligible condition when available. Identify whether assessment information was previously inspected; avoid relabeling retrospective cases as blind.
\item Show what the temperature comparison teaches about boundary conditions, loss paths and observation mapping. A CFD-conditioned map, measured-startup trajectory and autonomous forward prediction receive separate labels and scores.
\end{itemize}

\subsection{Developing profiles and conditional wall response}\label{subsec:developing-profiles-and-conditional-wall-response}
\begin{itemize}[leftmargin=*]
\item \researchlabel{developing-profile-and-transfer-studies}\textbf{Developing-profile and transfer studies.} Use temperature, transport, and buoyancy comparisons to assess whether a low-order profile representation warrants extension to the TRIDENT loops.
\begin{itemize}
\item TAMU geometry and fluid, inherited hydraulic and exterior-boundary laws, prescribed cooler inputs, and numerical-refinement requirements limit direct transfer to other facilities.
\end{itemize}

\item \researchlabel{conditional-wall-response}\textbf{Conditional wall and exterior-boundary studies.} Compare radial, regional, downcomer, and exterior-boundary representations using explicit fitted quantities, initialization scenarios, and regional errors.
\begin{itemize}
\item Measured primary temperatures are supplied inputs; insulated-region shell response does not resolve every bare-patch contact sensor.
\end{itemize}

\item VCU run-condition, boundary, regional-resistance, and downcomer comparisons provide further preliminary evidence. Their supplied temperatures and boundary assumptions must accompany any interpretation.
\item Use the isolated TAMU developing-profile work as a mechanism/\allowbreak{}reduction benchmark with explicit transfer limits.
\begin{itemize}
\item Distinguish temperatures associated with transported enthalpy, area-averaged density and buoyancy, and local probe response. Assess the thermal and hydraulic assumptions separately.
\end{itemize}
\end{itemize}
